% !TeX program = xelatex
\documentclass[margin = 5mm, tikz]{standalone}
\usepackage{xfrac}
\usepackage[MnSymbol]{mathspec}
\usepackage{esvect}
\setmainfont[Mapping=tex-text,Ligatures={NoRequired,NoCommon,NoContextual}]{Comic Sans MS}
\setmathfont(Digits,Latin,Greek)[Numbers={Lining,Proportional}]{Comic Sans MS}
\usetikzlibrary{shapes.geometric,decorations,decorations.pathmorphing}
\usetikzlibrary{intersections}
\usetikzlibrary{patterns}
\usetikzlibrary{calc}
\usetikzlibrary{shapes.misc}
\usetikzlibrary{spy}

\pgfdeclarelayer{bg}
\pgfdeclarelayer{bbg}
\pgfsetlayers{bbg, bg, main}

\definecolor{c1}{HTML}{ac0000}
\definecolor{c2}{HTML}{00a0d5}

\newcommand{\abs}[1]{{\left|{#1}\right|}}

\makeatletter
\newcommand\binomialCoefficient[2]{%
	% Store values 
	\c@pgf@counta=#1% n
	\c@pgf@countb=#2% k
	%
	% Take advantage of symmetry if k > n - k
	\c@pgf@countc=\c@pgf@counta%
	\advance\c@pgf@countc by-\c@pgf@countb%
	\ifnum\c@pgf@countb>\c@pgf@countc%
		\c@pgf@countb=\c@pgf@countc%
	\fi%
	%
	% Recursively compute the coefficients
	\c@pgf@countc=1% will hold the result
	\c@pgf@countd=0% counter
	\pgfmathloop% c -> c*(n-i)/(i+1) for i=0,...,k-1
	\ifnum\c@pgf@countd<\c@pgf@countb%
		\multiply\c@pgf@countc by\c@pgf@counta%
		\advance\c@pgf@counta by-1%
		\advance\c@pgf@countd by1%
		\divide\c@pgf@countc by\c@pgf@countd%
	\repeatpgfmathloop%
	\the\c@pgf@countc%
}
\makeatother

\tikzset{
    point/.style n args = {1}{circle, draw = black, inner sep = #1pt, fill = black},
    empty point/.style = {circle, draw = black, inner sep = 2pt, fill = white},
    cross/.style = {cross out, draw, minimum size = 2 * (#1 - \pgflinewidth), inner sep = 0pt, outer sep = 0pt},
    xkcd/.style n args = {1}{decorate, decoration = {random steps, pre length = #1mm, post length = #1mm, segment length = 0.6mm, amplitude = 0.2pt}},
    xkcd/.default = {4},
    point/.default = {2}
}

\begin{document}
    \begin{tikzpicture}
    	\pgfmathsetmacro{\dx}{1.7}
    	\pgfmathsetmacro{\dy}{1}
    	\foreach \i in {0,...,11} {
    		\foreach \j in {0,...,\i} {
    			\pgfmathsetmacro{\xc}{(\j - 0.5 * \i) * \dx}
    			\pgfmathsetmacro{\yc}{-\i * \dy}
    			\node (\i\j) at (\xc, \yc) {$\binomialCoefficient{\i}{\j}$};
	    	}
    	}
    	\foreach \i in {2,...,11} {
    		\pgfmathtruncatemacro{\upper}{\i - 1}
    		\foreach \j in {1,...,\upper} {
    			\pgfmathtruncatemacro{\upleft}{\j - 1}
    			\pgfmathtruncatemacro{\upright}{\j}
    			\draw[xkcd, thick, ->, c1] (\upper\upleft) -- (\i\j) node[circle, draw=c1, inner sep = 0.01cm, fill=white, pos=0.5]{\small$+$};
    			\draw[xkcd, thick, ->, c2] (\upper\upright) -- (\i\j) node[circle, draw=c2, inner sep = 0.01cm, fill=white, pos=0.5]{\small$+$};
    		}
    	}
    \end{tikzpicture}
\end{document}